\documentclass[11pt]{article}
\usepackage[a4paper,margin=25mm]{geometry}
\usepackage{amsmath,amssymb,amsthm,booktabs,longtable,hyperref,microtype}
\setlength{\emergencystretch}{3em}
\newtheorem{theorem}{Theorem}
\newtheorem{lemma}[theorem]{Lemma}
\newtheorem{proposition}[theorem]{Proposition}
\newtheorem{remark}[theorem]{Remark}
\newcommand{\Q}{\mathbb Q}\newcommand{\Z}{\mathbb Z}
\title{A shared-reference quartic checkpoint\\with an explicit source-universality obligation}
\author{\(\mathbb Z\mathbb Q\) Diophantine Project}
\date{23 November 2026}
\begin{document}\maketitle
\begin{abstract}
We construct complete rational quadratic systems with 2086, 2278, and
846 coordinates and, respectively, 1776, 1944, and 728 residuals. Their
sums of squares are quartics. The construction shares the actual elliptic
reference point among twelve integer-certification blocks, after proving
a simultaneous witness theorem under the source arithmetic interfaces.
The last route additionally assumes ordinary rational \(abc\).

The undecidability interpretation is conditional on a separate, explicit
hypothesis \(\mathcal U_{\rm em}\): universality of the precise emended
positive-integer source system. The consulted displays disagree in a
coefficient and contain a separate Pell-row typo. That source-verification
obligation has not been discharged. The numerical counts and finite
algebraic transformations are verified independently of
\(\mathcal U_{\rm em}\); this paper does not claim an unconditional
undecidability bound or a global minimum.
\end{abstract}

\section{Scope and the exact unresolved source assertion}
The independent-block construction is recorded in the companion paper
\nolinkurl{jones1980_quartic_transfer.tex} and its complete tables. Its
optimized outer circuit has 251 coordinates and 219 quadratic residuals,
with eleven ordinary integer inputs and one reciprocal integer input.
The present paper changes the way those twelve integrality requirements
are implemented.

There are two distinct kinds of assumptions. The integer-defining
interfaces retain their stated elliptic, height, finite prime-pattern,
recurrence, and native-pole premises; the 56-coordinate interface also
uses ordinary rational \(abc\). In addition, an undecidability conclusion
from this particular outer circuit requires the following assertion.
\begin{quote}
\textbf{Hypothesis \(\mathcal U_{\rm em}\).}
The precise displayed 28-variable positive-integer system used in the
companion paper, with
\[
 d=bw+ca-2c+(4a-5)\gamma,
 \qquad (p^2-1)k^2+1=\tau^2,
\]
has the asserted universality property: for some fixed positive integer
coefficients \(z,u,y\), the positive input values \(x\) for which it is
soluble form a nonrecursive set.
\end{quote}
Jones's 1980 announcement prints Greek \(\alpha\) in place of Latin
\(a\) in the coefficient of \(4\gamma\). Matiyasevich's 1999 survey
supports the Latin-\(a\) form but prints a separate, impossible equality
in the Pell row. No complete consulted display is therefore literally
identical to the emended system used here. The companion source audit
states the discrepancy and the remaining proof obligation in full.
The local Pell identities supporting the emendations do not themselves
prove \(\mathcal U_{\rm em}\).

Throughout, an undecidability statement about the final quartics is to
be read with \(\mathcal U_{\rm em}\) and the selected integrality
interface's premises. The simultaneous integer-witness theorems below
concern the integrality interfaces themselves and do not need
\(\mathcal U_{\rm em}\). The composition with the emended outer circuit
is an existence equivalence; its undecidability label additionally needs
that source hypothesis. Counts include every existential rational
unknown. The instance parameter \(x\), restricted to the cofinite range
\(x>2z\), is a coefficient parameter and is not an unknown.

\paragraph{Attribution.}
The defining blocks build on the supplied notes of Gabe Conant with
Astra 6.2 and the subsequent project constructions. This paper proves
new simultaneous-witness and register-sharing deductions from their
explicitly retained interfaces; it does not replace their arithmetic
hypotheses by finite computation.

\section{Simultaneous integral witnesses with one reference point}
\label{sec:shared-reference}

The general 163-coordinate defining system admits a genuine simultaneous
improvement.  Its reference point need not be chosen independently for
each input.  The assertion below uses the arithmetic interfaces retained
in the source manuscript
\texttt{scalar\_point163\_179\_20261120\_v82.tex}; it does not establish
those interfaces.  In particular, its prime supplier is required to
apply to every fixed finite family of rational contact positions, with
arbitrary finite avoidance and the compatible annihilation conditions.
No uniformity in a family that has not yet been fixed is needed.

\begin{theorem}[Shared reference for finitely many inputs]
\label{thm:shared-reference}
Fix a positive integer $k$ and a parameter $N$ for which the retained
general 163-coordinate system defines $\mathbb Z$ over $\mathbb Q$.
In $k$ copies of that system, identify the ten stored reference
coordinates
\[
                 15,16,17,18,19,20,50,51,103,104.
\]
The resulting quadratic system has rational input projection
$\mathbb Z^k$.  It can be written with
\[
                  153k+10\quad\hbox{coordinates},\qquad
                  129k+8\quad\hbox{quadratic residuals}.
\]
The coordinate count includes all $k$ distinguished inputs.  The sum
of the squares of these residuals is a quartic with the same input
projection.  The retained parameter threshold is unchanged.
\end{theorem}

\paragraph{Which coordinates are identified.}
The physical source labels of these coordinates are, respectively,
\[
                 15,16,17,18,19,20,64,65,150,151.
\]
In the literal map in
\nolinkurl{computation/variable_reductions_v82/general/quartic163/maps.json},
every one has $c$-weight $1$, $R$-weight $0$, and parameter-weight $0$.
The physical reference functions are source coordinates divided by
$c$.  Their stored coordinates are therefore the actual reference
functions, independent of the input-dependent scalar normalizer, of
the contact parameter, and of the contact gauge $c$.  Six coordinates
give the retained reference chart, two give its reciprocal square,
and two give the alternate reference ratio.  Once the actual point
$H$ is fixed, all ten have fixed values; this is not an identification
of freely chosen elliptic points.

The proof must verify a common choice of $H$, rather than merely
observe that these coordinates have the same names.  We give that
verification in the order in which the choices are made.

\begin{lemma}[Joint scalar preparation]
\label{lem:joint-scalar-preparation}
For any finite list of integer inputs $x_1,\ldots,x_k$, the ordered
scalar preparation in the 163-coordinate construction can be carried
out with one reference point $H=[h]P_0$ and one scalar difference
point $Q_2=[a]R_0$, with $a\in6\mathbb Z_{>0}$.  Every input retains
its own normalizer $R_i$, scalar quotient $\tau_i$, and scalar pole
witnesses.  The integer $a$ may be arbitrarily large.
\end{lemma}
\begin{proof}
Write
\[
 \beta_i=51x_i^2+2504,\qquad
 r_i=\beta_i(\beta_i-1)(\beta_i-2401)(32\beta_i+343),\qquad
 R_i=\frac{149^{2N+1}r_i}{1+149^2r_i}.
\]
All these numbers are fixed before any comparison prime is selected.
For $N\ge1$ the source estimate gives $R_i>100$.
Let $\mathcal P$ be the finite union, over all $i$, of the primes
where $R_i(R_i+9792)$ has odd valuation, together with $2,3,17$.
Every odd valuation in that union is positive.  Moreover every
$R_i(R_i+9792)$ is a square at $2,3,17$, by the same integer-input
congruences as in the one-input proof.

Choose $H$ sufficiently deep at both places over every $p\in\mathcal
P$, with $64\mid h$, and within the source real neighborhood of the
identity.  Write $d_w=v_w(z(H))$ at a place $w$ over $p$.  For each
$p$, choose one positive exact depth
\[
                         t_p=v_p(a/h).
\]
The lower bound on $t_p$ is the maximum of the finitely many lower
bounds needed for the inputs $i=1,\ldots,k$.  The reference depths
are chosen first sufficiently large that this lower bound is smaller
than
\[
             2\min_{w\mid p} d_w+
             \tfrac12\min_i v_p(R_i)-q/2,
\]
where $q\ge2$ is any fixed required accuracy.  Such a choice is
possible since the lower bounds are fixed before the reference depths
are increased.  This gives, simultaneously for every $i$ and every
$w\mid p$,
\[
 v_w(\epsilon_w)\ge q,\qquad
 v_p(R_i)+4d_w-2t_p\ge q,
 \qquad Z=(a/h)(1+\epsilon_w).
\]
These are exactly the two inequalities used in the source scalar
preparation, now with their maxima and minima taken over a finite
list.  Exact depth is a nonempty open congruence condition.  The
Chinese remainder theorem combines these conditions with $a=2a_0$,
$3\mid a_0$, and the fixed component and reduction-order congruences.
The retained recurrence gives arbitrarily large positive choices of
$a_0$ in that progression with $Q_2=[a_0]P_0$ in the source real
neighborhood.

The scalar trace calculation consequently gives, for every $i$,
\[
              v_p(\tau_i)=2t_p-5v_p(R_i)-v_p(2).
\]
The common lower bound on $t_p$ makes all $\tau_i$ integral on the
odd supports and gives the prescribed depths at $2,3,17$.  Outside
$\mathcal P$ the parity argument for the reduced denominator applies
separately to each $R_i(R_i+9792)\tau_i^2$.  There is no requirement
that the different $\tau_i$ be equal.

The source real neighborhoods can also be common: they give
$|f(H)^{(j)}|\le1$ and $1\le|Z^{(j)}|\le3$ at both real embeddings,
and hence for every $i$
\[
       0<R_i(R_i+9792)\tau_i^2
         \le\frac{24482500}{R_i^4}<17.
\]
Thus every actual scalar input satisfies its realization hypotheses.
After $a$ is fixed these conditions are open in $H$.  Intersect the
finitely many reference neighborhoods and retain a progression through
the chosen $h$ preserving all of them.  This intersection is nonempty
because it contains that chosen reference point.  It is fixed before
comparison primes are selected.
\end{proof}

\begin{lemma}[Separating all fixed contact lists]
\label{lem:contact-offsets}
Let $\Theta\subset\mathbb Q$ be the source contact set, of diameter
four, and suppose $T_1,\ldots,T_k>7$ are fixed rational numbers.
There are integers $\delta_1,\ldots,\delta_k$ such that the $2k$
sets
\[
       \delta_i+\Theta,\qquad \delta_i+T_i+\Theta
                      \quad(1\le i\le k)
\]
are pairwise disjoint.  Consequently the retained fixed-family prime
supplier applies once to their union.
\end{lemma}
\begin{proof}
The two sets belonging to a single $i$ are disjoint since $T_i$ is
larger than the diameter.  Having chosen the first $i-1$ integers,
choose $\delta_i$ so large that both new sets lie strictly to the
right of every previously chosen set.  This proves the assertion
by induction.  All positions are fixed rational numbers before the
prime supplier is invoked.
\end{proof}

\begin{proof}[Proof of Theorem~\ref{thm:shared-reference}]
Soundness is immediate from the one-input theorem: every retained
block still satisfies every original residual, with the identified
coordinates simply assigned common values.  Thus each of its inputs
is integral.  This argument does not require distinct inputs or
distinct contact radicals on an arbitrary rational solution.

For completeness fix an arbitrary tuple in $\mathbb Z^k$.  Apply
Lemma~\ref{lem:joint-scalar-preparation}.  The common scalar label
$a\in6\mathbb Z_{>0}$ then supplies the common Lucas labels $t$ and
odd $b>1$ from the source recurrence.  Put
\[
 Q_{1,i}=[\beta_i]P_0,\qquad Q_2=[a]R_0,\qquad
 Q_3=[t]R_0,\qquad Q_4=[(b-1)/2]R_0.
\]
These points are fixed and nonzero nontorsion.  Every input retains
its own aggregate
\[
 T_i=7+\frac14\operatorname{Tr}
          \bigl(x(Q_{1,i})+R_i x(Q_2)+x(Q_3)+x(Q_4)\bigr)>7.
\]
All forbidden finite chart points can be avoided by shrinking the
common reference neighborhood.  There are only finitely many such
conditions; multiplication by each fixed nonzero root label is
nonconstant on the relevant connected real orbit closure.

Keep the contact parameter $\lambda_i$ of every block independent
in the equations.  Choose the integers $\delta_i$ in
Lemma~\ref{lem:contact-offsets} and, solely when constructing this
witness, set $\lambda_i=\lambda+\delta_i$.  The $\delta_i$ are
not additional quantified variables or extra equations.  The single
finite-family prime supplier is applied to the union of these fixed
translated positions, after enlarging finite avoidance to include
all the normalizers, all fixed nonzero root multipliers, all fixed
chart denominators, and the previous preparation moduli.

The successful reference annihilation conditions concern $h$ and
$h-1$.  They are the same conditions for every input: the primary
reference is formal in one slot and its translate by $P_0$ is formal
in the other slot.  The different $\beta_i$ multiply the same
generator and are units at the selected primes; they prescribe no
conflicting residues for $h$.  The source half-generator conditions
are retained for the common points $Q_2,Q_3,Q_4$.

The source short-interval argument makes all selected primes
coprime to every selected reduction order and to the earlier
preparation moduli.  The same synchronization therefore applies to
the entire finite family.  At each selected prime only the smaller
of the two initial reference depths is increased, by specifying the
valuation of $h$ or of $h-1$; equal initial depths require neither
to be divisible by the prime.  These additional prime-power
conditions are mutually coprime and preserve every earlier
annihilation congruence.  The source compact-component recurrence
preserves the real return on the resulting progression.

For each block, its two contact lists now have the same properties
used in its original proof.  At a selected prime belonging to that
block, its successful and opposite reference depths are synchronized,
all required root multipliers are units, and all its detector terms
have the original positive depth.  Thus its comparison inputs have
selected valuation one.  At fixed places the expansion
\[
 q(\lambda+u)=\kappa\lambda^{-2}+O(\lambda^{-3}),\qquad \kappa>0,
\]
holds for every one of the finitely many fixed shifts $u$; the source
local preparation and real limits apply to all blocks simultaneously.
Every actual scalar and comparison-pole input receives its own native
realization witnesses.  Filling the complete finite charts from the
one final $H$ gives identical values in the ten reference coordinates.
This proves completeness, including tuples with repeated entries.

Identifying ten coordinates in all $k$ copies saves $10(k-1)$,
so the total is $163k-10(k-1)=153k+10$.  The six rows named
\texttt{sharedT.reference.relation}$j$.\texttt{$\ell$}, with
$1\le j\le3$ and $0\le\ell\le1$, become identical in all copies.
So do the two rows
\texttt{testSquare.40.0.0.definition} and
\texttt{testSquare.40.0.1.definition}.  Dropping their duplicate
copies leaves $137k-8(k-1)=129k+8$ residuals.  Every retained
residual is still quadratic.  The input-square residual has a
nonzero quadratic part, so their sum of squares has degree exactly
four over the real field $\mathbb Q$.
\end{proof}

\begin{remark}[Use inside a universal quadratic system]
If a quadratic outer system has $b$ rational coordinates, including
$k$ coordinates whose integrality is certified by unmodified copies
of the 163-coordinate system, the shared-reference construction has
\[
                    b+152k+10
\]
coordinates, compared with $b+162k$ for independent certificates.
This formula presupposes that no subsequent homogenization changes
the ten reference coordinates in only some blocks.  Such a block
can be excluded from the sharing, or its specific coordinate changes
must be accounted for separately.  Identifying coordinates merely
because their source labels agree after different gauges would not
justify this formula.
\end{remark}

\begin{remark}[What is not shared]
The argument does not identify the normalizers $R_i$, scalar pole
quotients, contact parameters, or contact gauges.  Independent contact
parameters are essential to the proof: setting every $\lambda_i$
equal would produce repeated contact poles when inputs coincide.
The offset construction uses their existing freedom to give a valid
simultaneous witness.  More extensive sharing of the scalar and
comparison circuits requires a separate joint compilation and cannot
be inferred from the ten-coordinate identification alone.
\end{remark}

\section{The fixed and ordinary-\(abc\) interfaces}

\subsection{The fixed 179-coordinate interface}
The shared-reference proof also applies to the fixed 179-coordinate
interface. Here the identified stored coordinates are
\[
                15,16,17,18,19,20,50,51,104,105.
\]
Their physical labels are again \(15,16,17,18,19,20,64,65,150,151\),
all with \(c\)-weight one and \(R\)-weight zero. The same eight reference
rows are duplicated. The three-register rechart used in the combined
compiler changes none of these coordinates or rows.

For completeness, the fixed normalizers are
\[
 R_i=A_*r_i/(1+25r_i)>100,
\]
with \(A_*\) the fixed source integer. Their square-class preparation is
identical to the one-input preparation: \(R_i(R_i+9792)\) has only positive
odd valuations, and is a square in \(\Q_2,\Q_3,\Q_{17}\). Take the finite
union of their odd supports, the maximum of all required lower scalar
depths, and reference formal depths large enough for all the resulting
strict inequalities. This supplies a single reference point and a single
positive scalar label for every input. The same common label supplies the
Lucas labels. Each block keeps its own scalar quotient and native pole
witnesses.

The effective construction has its original unweighted aggregate and its
original contact lists. After all point differences have been chosen,
these lists are fixed, finite, and internally separated by the source
construction. Translate the whole contact list of block \(i\) by a fixed
integer \(\delta_i\), choosing these integers successively so that different
blocks' lists are disjoint. This is achieved by setting that block's free
contact parameter equal to the common slope plus \(\delta_i\) solely in
the completeness construction. The fixed-family prime supplier applies
to the union. The shared reference annihilation conditions are the same
conditions on \(h\) and \(h-1\) as before, and finite avoidance is enlarged
to the union of the finitely many earlier exceptional sets. All fixed-place
and real preparations persist because their slope asymptotics are unchanged
by fixed translations. This proves simultaneous completeness. Soundness
uses each original 179-coordinate block separately, and hence retains its
effective coefficients and height bounds without alteration.

Consequently \(k\) copies have a simultaneous presentation in
\[
                       169k+10
\]
coordinates, including their \(k\) inputs, with \(143k+8\) quadratic rows.

\subsection{Nonzero inputs for the 56-coordinate ordinary-\(abc\) interface}
For this interface it suffices to prove simultaneous completeness for
nonzero integer inputs. This restricted assertion avoids the special
zero-input reference in the source; no assertion that every tuple involving
zero has a common reference is needed below.

\begin{proposition}\label{prop:abc-common-reference}
Fix a defining odd coefficient parameter of the 56-coordinate interface,
under ordinary rational \(abc\) and all its retained arithmetic premises.
Identify coordinates \(2,3,4,5,6,7\) in \(k\) copies of its quadratic
system. Every rational solution has integral inputs. Every tuple in
\((\Z\setminus\{0\})^k\) has a rational solution of this shared system.
It has \(50k+6\) coordinates, including the inputs, and \(42k+4\) rows.
\end{proposition}

\begin{proof}
The four rows named \texttt{reference.primitive.1} through
\texttt{reference.primitive.4} involve exactly coordinates \(2,\ldots,7\)
and fixed coefficients. These are the six actual reference-chart
coordinates; they do not depend on the input or the contact gauge. Each
original block remains a solution of its own integer-defining system,
so soundness is unchanged.

For completeness fix nonzero integers \(n_1,\ldots,n_k\). The parent
multiplication family uses
\[
 U_i=[n_i]H,\qquad V_i=[n_i](H-R_0),\qquad U_i-V_i=[n_i]R_0.
\]
Thus the scaled trace
\[
 T_i=n_i^2\operatorname{Tr}x([n_i]R_0)
\]
is fixed independently of the later reference choice. So are
\[
 M_i=B+2(n_i^8+T_i^2),\qquad a_i=M_i^4+16>16.
\]
The two-adic source lemma gives \(a_i\in\Q_2^{\times2}\) for every \(i\).

First choose a permissible common reference neighborhood. For a nonzero
label \(n_i\), multiplication by \(n_i\) on the elliptic curve is
nonconstant and has finite fibers. The forbidden finite chart points for
all the \(U_i,V_i,H,H-R_0\) therefore exclude only finitely many points
from the reference orbit. Choose one reference in the cyclic orbit avoiding
their finite union, and take sufficiently small real and fixed-place
neighborhoods in which all the finitely many raw tests are defined and
bounded. These are the bounded neighborhoods required by the source
prime-pattern construction. They have a common point by their construction;
an arbitrary intersection of independently chosen neighborhoods is not
being assumed nonempty.

For block \(i\), the sixteen fixed contact positions are
\[
       \pm(a_i+r),\qquad r\in\{-8,-7,-4,-1,1,4,7,8\}.
\]
They are internally distinct because \(a_i>16\). Choose integers
\(\delta_i\) so that the translated lists
\(\delta_i\pm(a_i+r)\) are pairwise disjoint over all blocks. Retain
independent contact variables \(z_i\) in the equations, and choose
\(z_i=z+\delta_i\) only for this witness. Enlarge the fixed preparation
set to include the supports of every \(n_i\), every contact position and
its nonzero pairwise differences, every \(a_i\), and the source chart
exclusions. Choose a common integer \(D_0>0\) clearing the finitely many
rational contact denominators. The retained fixed-family supplier then
applies to the \(16k\) distinct forms
\[
       N+D_0^2\bigl(\delta_i\pm(a_i+r)\bigr)L^2,
       \qquad z=N/(D_0^2L^2).
\]
All listed primes have the required finite avoidance. The reference
conditions at them are the same successful primary and opposite
conditions on \(h\) and \(h-1\). Since every \(n_i\) is a unit at the
selected primes, multiplication by \(n_i\) introduces no conflicting
reference congruence. The compatible annihilation refinements and the
compact-component recurrence are applied to this one finite union.

At each selected contact of block \(i\), exactly one of its sixteen
factors has simple vanishing. The original derivative identity therefore
gives contact-quotient valuation one, and the original point preparations
give positive valuation of that block's shifted norm/trace gate. All its
selected native-pole valuations are consequently the same as in the
one-input proof. The factorization argument outside its selected support
is unchanged; a prime selected for a different block is simply an
unselected prime in this argument.

At the finitely many fixed places, take the denominator of the common
slope \(z\) sufficiently deep for all blocks. Since \(\delta_i\) is fixed,
\(z_i/z\) tends to one at those places and at real infinity. Thus the
source asymptotics
\[
                 u_i\sim z_i^2/a_i,\qquad c_i\sim a_i/z_i^2
\]
hold simultaneously. The bounded raw point tests give
\(W_i-1=O(c_i)\), \(D_i=O(c_i^2)\), and the native-pole value is
asymptotic to \(7c_i\). At two, every \(a_i\) is a square, so the native
input lies in \(7\Q_2^{\times2}\) with the required arbitrarily large
positive valuation. At real infinity it lies in \((0,1)\), and the source
parity argument gives a square reduced denominator for each block.
Apply the stated native-pole realization separately to all these actual
inputs. The one final reference point gives the same six reference
coordinates in every block, proving completeness.

Identifying six coordinates saves \(6(k-1)\); deleting the four repeated
reference rows saves \(4(k-1)\). This gives the displayed counts.
\end{proof}

\section{Common references in the optimized Jones system}
The optimized independent-block construction has 251 outer coordinates and
219 outer rows. Its integer requirements consist of eleven ordinary inputs
and the reciprocal input \(1/v\). To share all references in the same
coordinate scale, introduce one further coordinate \(i\), impose
\[
                              iv=1,
\]
and use an ordinary integer block for \(i\). This replaces the homogeneous
reciprocal block, whose witness coordinates were scaled by \(v\). No
additional arithmetic assumption is introduced: the new equation makes
\(i=1/v\) exactly.

All twelve integer blocks are now unhomogenized and may share their actual
reference coordinates. For the 163- and 179-coordinate interfaces, identify
the ten coordinates just proved to be simultaneous. For the 56-coordinate
interface, use Proposition~\ref{prop:abc-common-reference}. Its nonzero-input
restriction is sufficient: in the completeness direction choose a positive
integer witness of Jones's original universal system, before applying the
sign and reciprocal changes. All twelve inputs to the integer blocks are
then positive. In the soundness direction, every individual block still
certifies integrality, including if a relaxed outer solution initially has
some zero input. Thus no unproved common-reference statement at zero is
used.

The resulting full counts are
\begin{center}
\begin{tabular}{rrrrr}
\toprule
Interface&Outer coordinates&Shared coordinates saved&Total coordinates&Total rows\\
\midrule
163&252&110&2086&1776\\
179&252&110&2278&1944\\
56&252&66&846&728\\
\bottomrule
\end{tabular}
\end{center}
For example the first count is \(252+12(163-1)-110=2086\), while its row
count is \(220+12\cdot137-88=1776\). The analogous ordinary-\(abc\) count
is \(252+12(56-1)-66=846\), with \(220+12\cdot46-44=728\) rows.

The separate compiler \texttt{build\_shared\_reference.py} starts from the
complete independent-block tables, replaces the reciprocal block by its
ordinary version and the literal row \(iv=1\), identifies precisely the
proved reference coordinates, and verifies that every discarded duplicate
row is identical to its retained representative. It leaves the original
tables unchanged. All residuals remain ordinary quadrics; all coordinates
occur. Their sum of squares has degree exactly four and the required
existence equivalence. The source-universality hypothesis \(\mathcal U_{\rm em}\), the source arithmetic premises, the ordinary
\(abc\) hypothesis for the 56-coordinate interface, and the cofinite
restriction on the instance parameter are all retained.

\section{Reproduction and the limits of the certificate}
From the repository root, run
\begin{verbatim}
python -B undecidability/verify_shared_reference.py
\end{verbatim}
The verifier reconstructs every table from the preserved independent-block
baseline, checks the extra literal reciprocal equation, checks every
discarded row against its retained representative, and independently
inspects integer coefficient polynomials, quadratic degrees, coordinate
support, and the constant quadratic terms ensuring exact degree four.
It also verifies that the source-universality qualification is present
in every table and receipt. A PASS certifies these finite facts. It does
not certify \(\mathcal U_{\rm em}\), the retained infinite arithmetic
interfaces, or minimality.

\begin{thebibliography}{9}
\bibitem{Jones80} J. P. Jones, \emph{Undecidable Diophantine equations},
Bulletin of the American Mathematical Society (New Series)
\textbf{3} (1980), 859--862.
\href{https://doi.org/10.1090/S0273-0979-1980-14832-6}{doi:10.1090/S0273-0979-1980-14832-6}.
\bibitem{Jones82} J. P. Jones, \emph{Universal diophantine equation},
Journal of Symbolic Logic \textbf{47} (1982), 549--571.
\href{https://doi.org/10.2307/2273588}{doi:10.2307/2273588}.
\bibitem{Mat99} Yu. Matiyasevich, \emph{Hilbert's tenth problem: A two-way
bridge between number theory and computer science}, in C. S. Calude (ed.),
\emph{People and Ideas in Theoretical Computer Science}, Springer,
Singapore, 1999, pp. 177--204.
\url{https://www.cs.auckland.ac.nz/~cristian/tcspi/ch9.pdf}.
\bibitem{General} \(\mathbb Z\mathbb Q\) Diophantine Project,
\emph{A scalar-point sharing construction in 163 and 179 variables},
20 November 2026, repository manuscript
\nolinkurl{papers/quartic/scalar_point163_179_20261120_v82.tex}.
\bibitem{ABC} \(\mathbb Z\mathbb Q\) Diophantine Project,
\emph{Paired contact coordinates: an ordinary-abc quartic in 56 variables},
19 November 2026, repository manuscript
\nolinkurl{papers/quartic_existence/paired_contact56_20261119_v81.tex}.
\bibitem{Companion} \(\mathbb Z\mathbb Q\) Diophantine Project,
\emph{A direct quadratic compilation of the emended universal system},
23 November 2026, companion checkpoint and source reconciliation,
\nolinkurl{undecidability/papers/jones1980_quartic_transfer.tex}.
\end{thebibliography}
\end{document}
